% !TeX root = ../../midterm-report.tex
\subsection{实验结果与讨论}
表~\ref{tab:example} 展示三线表，图~\ref{fig:example} 展示双图横排占位。请用真实结果替换，不将占位内容当成数据。
\begin{table}[htbp]
\centering\wuhao
\caption{结果表格占位示例}\label{tab:example}
\begin{tabular}{lll}
\toprule[1.5pt]
方法 & 指标一 & 指标二\\
\midrule[1pt]
基线方法 & 待填写 & 待填写\\
待比较方法 & 待填写 & 待填写\\
\bottomrule[1.5pt]
\end{tabular}
\end{table}
表格后说明比较条件与结论来源；必要时报告重复次数、统计口径及不确定性，不单独贴表。

\begin{figure}[htbp]
\centering
% 自包含占位框，无需迁入原报告图片；替换时用 0.45\linewidth 左右的图片。
\subfigure[子图一]{\fbox{\parbox[c][1.5cm][c]{0.40\linewidth}{\centering 图像占位}}}
\hfill
\subfigure[子图二]{\fbox{\parbox[c][1.5cm][c]{0.40\linewidth}{\centering 图像占位}}}
\caption{横排子图占位示例}\label{fig:example}
\end{figure}
图中仅提供版式位置。真实图表应有必要的前文引入与后文解释；优先横排、裁除无效边距或拆分过大的综合图，不滥用强制浮动和换页。
